%! Unit 1: Linear Functions — Unit Test (KEY)
% GENERATED BODY: the blank and key bodies are byte-identical except the header title;
% answers live in \slot / \optok / work, which the preamble shows or hides.
\documentclass[12pt]{article}
\usepackage{algebra2-article}
\usepackage{algebra2-key}   % pulls in -boxes; do NOT also load -boxes
% Coordinate grid: {xmin}{xmax}{ymin}{ymax}
\newcommand{\lgrid}[4]{%
  \draw[step=1, linegray!40, very thin] (#1,#3) grid (#2,#4);
  \draw[->, semithick, charcoal] (#1,0)--(#2,0) node[below right, font=\scriptsize]{$x$};
  \draw[->, semithick, charcoal] (0,#3)--(0,#4) node[left, font=\scriptsize]{$y$};}
% Endpoint dots: {x}{y}. Closed = filled; open = white-filled ring.
\newcommand{\fdot}[2]{\fill[forest] (#1,#2) circle (4.5pt);}
\newcommand{\hdot}[2]{\draw[very thick, forest, fill=white] (#1,#2) circle (4.5pt);}
% Part-divider strip (headlinebox takes one arg: the background color)
\newcommand{\parthead}[1]{%
  \boxguard[10]\vspace{6pt}\noindent
  \begin{headlinebox}{forest}\color{white}\bfseries #1\end{headlinebox}%
  \vspace{4pt}}
% THE WORK RULE for a test: \slot{width}{answer} and \opt / \optok are the ONLY
% lines that differ between blank and key, and they differ in the preamble, not the body.
% A slot is a fixed-width underline in both files, so no line rewraps in the key.
\newcommand{\slot}[2]{\underline{\makebox[#1]{\ans{#2}}}}
\newcommand{\opt}[2]{(#1)~#2}
\newcommand{\optok}[2]{\mbox{\llap{\textcolor{keyred}{$\boldsymbol{\rightarrow}$}\,}\textcolor{keyred}{\textbf{(#1)}~#2}}}

\begin{document}

\pageheader{Unit 1: Linear Functions}{Unit Test --- Answer Key}
\namedateperiod

\vspace{0.06in}
\noindent\textbf{Instructions:} Show all work. Write every line in $y=mx+b$ form when asked.
No notes or calculator unless your teacher says so. \hfill\textbf{Total: 100 pts}

% ======================================================================
\parthead{Part A --- Types of Slope (8 pts)}
For each line, write whether its slope is \textbf{positive}, \textbf{negative}, \textbf{zero}, or
\textbf{undefined}.

\vspace{6pt}
\noindent
\begin{minipage}[t]{0.24\textwidth}\centering
\begin{tikzpicture}[scale=0.34]
  \lgrid{-3}{3}{-3}{3}
  \draw[very thick, navy] (-3,-2)--(2.5,3);
\end{tikzpicture}\\[2pt]
1.\ \slot{2.4cm}{positive}
\end{minipage}%
\begin{minipage}[t]{0.24\textwidth}\centering
\begin{tikzpicture}[scale=0.34]
  \lgrid{-3}{3}{-3}{3}
  \draw[very thick, navy] (-3,2)--(3,2);
\end{tikzpicture}\\[2pt]
2.\ \slot{2.4cm}{zero}
\end{minipage}%
\begin{minipage}[t]{0.24\textwidth}\centering
\begin{tikzpicture}[scale=0.34]
  \lgrid{-3}{3}{-3}{3}
  \draw[very thick, navy] (-1.5,-3)--(-1.5,3);
\end{tikzpicture}\\[2pt]
3.\ \slot{2.4cm}{undefined}
\end{minipage}%
\begin{minipage}[t]{0.24\textwidth}\centering
\begin{tikzpicture}[scale=0.34]
  \lgrid{-3}{3}{-3}{3}
  \draw[very thick, navy] (-2,3)--(1,-3);
\end{tikzpicture}\\[2pt]
4.\ \slot{2.4cm}{negative}
\end{minipage}

\vspace{10pt}
\noindent No graph this time --- decide from the points or the equation.
\begin{enumerate}[leftmargin=1.9em, itemsep=6pt, label=\arabic*., start=5]
  \item The line through $(1,\,4)$ and $(1,\,-3)$: \hfill \slot{2.4cm}{undefined}
  \item The line $y=-2x+7$: \hfill \slot{2.4cm}{negative}
  \item The line $y=6$: \hfill \slot{2.4cm}{zero}
  \item The line through $(-2,\,1)$ and $(3,\,5)$: \hfill \slot{2.4cm}{positive}
\end{enumerate}

% ======================================================================
\parthead{Part B --- Writing Equations of Lines (20 pts)}
Write each line in slope-intercept form, $y=mx+b$. Show how you found $m$ and $b$.

\begin{enumerate}[leftmargin=1.9em, itemsep=10pt, label=\arabic*., start=9]
  \item The line through $(2,\,3)$ and $(4,\,11)$.
    \begin{work}
      m &= \frac{11-3}{4-2} = \frac{8}{2} = 4 \\
      3 &= 4(2) + b \quad\Rightarrow\quad b = -5
    \end{work}
    \vspace{0.5cm}
    \hfill Equation: \slot{3.6cm}{$y=4x-5$}

  \item The line through $(-2,\,7)$ and $(4,\,-2)$.
    \begin{work}
      m &= \frac{-2-7}{4-(-2)} = \frac{-9}{6} = -\frac{3}{2} \\
      7 &= -\frac{3}{2}(-2) + b = 3 + b \quad\Rightarrow\quad b = 4
    \end{work}
    \vspace{0.5cm}
    \hfill Equation: \slot{3.6cm}{$y=-\tfrac{3}{2}x+4$}

  \item The line with slope $3$ that passes through $(2,\,1)$.
    \begin{work}
      1 &= 3(2) + b \\
      1 &= 6 + b \quad\Rightarrow\quad b = -5
    \end{work}
    \vspace{0.5cm}
    \hfill Equation: \slot{3.6cm}{$y=3x-5$}

  \item The line with slope $-\tfrac{1}{2}$ that passes through $(4,\,6)$.
    \begin{work}
      6 &= -\frac{1}{2}(4) + b \\
      6 &= -2 + b \quad\Rightarrow\quad b = 8
    \end{work}
    \vspace{0.5cm}
    \hfill Equation: \slot{3.6cm}{$y=-\tfrac{1}{2}x+8$}
\end{enumerate}

% ======================================================================
\parthead{Part C --- Piecewise Functions (20 pts)}

\begin{enumerate}[leftmargin=1.9em, itemsep=10pt, label=\arabic*., start=13]
  \item \textbf{Evaluate.} (8 pts) \quad
    $f(x)=\begin{cases} -2x+1, & x<0 \\ 4, & 0\le x\le 2 \\ x+3, & x>2 \end{cases}$
    \\[8pt]
    $f(-3)=$ \slot{1.4cm}{$7$} \qquad $f(0)=$ \slot{1.4cm}{$4$} \qquad
    $f(2)=$ \slot{1.4cm}{$4$} \qquad $f(5)=$ \slot{1.4cm}{$8$}

  \item \textbf{Domain.} (12 pts) Give the domain of each function as an inequality or in interval
    notation.
    \begin{enumerate}[label=(\alph*), leftmargin=1.9em, itemsep=12pt]
      \item $g(x)=\begin{cases} x+1, & x<-2 \\ 3, & -2\le x<5 \end{cases}$
        \hfill Domain: \slot{3.8cm}{$x<5$, \ $(-\infty,\,5)$}
      \item $h(x)=\begin{cases} 2x, & 0\le x\le 3 \\ 9-x, & 3<x\le 7 \end{cases}$
        \hfill Domain: \slot{3.8cm}{$0\le x\le 7$, \ $[0,\,7]$}
      \boxguard[9]
      \item The function $k$ graphed below.
        \\[4pt]
        \begin{minipage}[c]{0.45\linewidth}\centering
        \begin{tikzpicture}[scale=0.5]
          \lgrid{-5}{4}{-3}{4}
          \draw[very thick, forest] (-4,-2)--(1,3);
          \fdot{-4}{-2} \hdot{1}{3}
          \draw[very thick, forest] (1,-1)--(3,-1);
          \fdot{1}{-1} \hdot{3}{-1}
        \end{tikzpicture}
        \end{minipage}%
        \hfill Domain: \slot{3.8cm}{$-4\le x<3$, \ $[-4,\,3)$}
    \end{enumerate}
\end{enumerate}

% ======================================================================
\parthead{Part D --- The Absolute Value Function (28 pts)}

\begin{enumerate}[leftmargin=1.9em, itemsep=10pt, label=\arabic*., start=15]
  \item (3 pts) Circle the letter of the graph of an \textbf{absolute value} function.
    \\[4pt]
    \begin{minipage}[t]{0.24\linewidth}\centering
    \begin{tikzpicture}[scale=0.32]
      \lgrid{-3}{3}{-3}{3}
      \draw[very thick, navy, domain=-2.5:2.5, smooth] plot (\x, {0.5*\x*\x-2});
    \end{tikzpicture}\\ \opt{A}{}
    \end{minipage}%
    \begin{minipage}[t]{0.24\linewidth}\centering
    \begin{tikzpicture}[scale=0.32]
      \lgrid{-3}{3}{-3}{3}
      \draw[very thick, navy] (-3,-2.5)--(3,1.5);
    \end{tikzpicture}\\ \opt{B}{}
    \end{minipage}%
    \begin{minipage}[t]{0.24\linewidth}\centering
    \begin{tikzpicture}[scale=0.32]
      \lgrid{-3}{3}{-3}{3}
      \draw[very thick, navy] (-2,1)--(1,-2)--(3,0);
    \end{tikzpicture}\\ \optok{C}{}
    \end{minipage}%
    \begin{minipage}[t]{0.24\linewidth}\centering
    \begin{tikzpicture}[scale=0.32]
      \lgrid{-3}{3}{-3}{3}
      \foreach \k in {-2,...,1} { \draw[very thick, navy] (\k,\k)--(\k+1,\k);
        \fill[navy] (\k,\k) circle (4pt); \draw[thick, navy, fill=white] (\k+1,\k) circle (4pt); }
    \end{tikzpicture}\\ \opt{D}{}
    \end{minipage}

  \item (3 pts) Which equation is an \textbf{absolute value} function?
    \\[4pt]
    \opt{A}{$y=3x-2$} \hfill \optok{B}{$y=|x-4|+1$} \hfill \opt{C}{$y=x^2+1$} \hfill
    \opt{D}{$y=\lfloor x\rfloor$}

  \item (15 pts) Each function is a transformation of the parent $y=|x|$. Fill in the table.
    \\[6pt]
    {\renewcommand{\arraystretch}{1.9}
    \begin{tabular}{|c|c|c|c|c|}
      \hline
      \textbf{Function} & \textbf{Vertex} & \textbf{Left/right shift} & \textbf{Up/down shift} &
      \textbf{Opens} \\ \hline
      $g(x)=|x-3|+2$ & \slot{1.9cm}{$(3,\,2)$} & \slot{2.1cm}{right 3} & \slot{2.1cm}{up 2} &
        \slot{1.4cm}{up} \\ \hline
      $h(x)=|x+4|-1$ & \slot{1.9cm}{$(-4,\,-1)$} & \slot{2.1cm}{left 4} & \slot{2.1cm}{down 1} &
        \slot{1.4cm}{up} \\ \hline
      $k(x)=-|x|+5$ & \slot{1.9cm}{$(0,\,5)$} & \slot{2.1cm}{none} & \slot{2.1cm}{up 5} &
        \slot{1.4cm}{down} \\ \hline
    \end{tabular}}

  \item (4 pts) Which equation matches the graph?
    \\[4pt]
    \begin{minipage}[c]{0.36\linewidth}\centering
    \begin{tikzpicture}[scale=0.36]
      \lgrid{-6}{2}{-2}{4}
      \draw[very thick, navy] (-5.5,-0.5)--(-2,3)--(1.5,-0.5);
      \fill[navy] (-2,3) circle (4pt);
    \end{tikzpicture}
    \end{minipage}%
    \begin{minipage}[c]{0.6\linewidth}
    \opt{A}{$y=-|x-2|+3$} \\[6pt]
    \optok{B}{$y=-|x+2|+3$} \\[6pt]
    \opt{C}{$y=|x+2|+3$} \\[6pt]
    \opt{D}{$y=-|x+3|+2$}
    \end{minipage}

  \item (3 pts) Compared with the parent $y=|x|$, the graph of $y=3|x|$ is
    \\[4pt]
    \optok{A}{narrower} \hfill \opt{B}{wider} \hfill \opt{C}{shifted up 3} \hfill
    \opt{D}{shifted right 3}
\end{enumerate}

% ======================================================================
\parthead{Part E --- Absolute Value Equations \& Inequalities (24 pts)}
Solve. Show both cases.

\begin{enumerate}[leftmargin=1.9em, itemsep=10pt, label=\arabic*., start=20]
  \item (8 pts) \quad $|2x-3|=7$
    \begin{work}
      2x-3 &= 7 \quad\text{or}\quad 2x-3 = -7 \\
      2x &= 10 \quad\text{or}\quad 2x = -4
    \end{work}
    \vspace{0.5cm}
    \hfill Solution: \slot{4.2cm}{$x=5$ or $x=-2$}

  \item (8 pts) \quad $|x+1|<4$
    \begin{work}
      -4 &< x+1 < 4 \\
      -5 &< x < 3
    \end{work}
    \vspace{0.5cm}
    \hfill Solution: \slot{4.2cm}{$-5<x<3$}

  \item (8 pts) \quad $|x-2|\ge 5$
    \begin{work}
      x-2 &\ge 5 \quad\text{or}\quad x-2 \le -5 \\
      x &\ge 7 \quad\text{or}\quad x \le -3
    \end{work}
    \vspace{0.5cm}
    \hfill Solution: \slot{4.2cm}{$x\ge 7$ or $x\le -3$}
\end{enumerate}

\end{document}
